\subsection{Closing the complete source translation}\label{sec:closed-program}
We can now connect the client calculation to the entire running program.
Use $x,f$ for its source binders and reserve $\ell,z$ for the two
locations allocated during execution. This keeps a parameter's assumed
interface distinct from a location's precise stored type.

\textbf{First supply the source object's interface.}
Let $D=\iface{W}{W}$. The source definition rule directly assigns
$A=W$ the exact declaration $D$. The object rule checks that definition
under a self binding at $\mu(s.D)$ that also records the literal's
definitions, and gives the literal type $\mu(s.D)$.
Here $W$ is independent of self, so checking the definition uses no
facts from that binding. After the first let, variable typing gives
$x:\mu(s.D)$, whereas reading a bound needs $x:D[x/s]$.
The source recursive-elimination rule supplies precisely that change:
\[
 \src{\frac{\Gamma\vdash x:\mu(s.D)}
                 {\Gamma\vdash x:D[x/s]}\ \textsc{Rec-E}.}
\]
It changes the derivation of the variable occurrence, not the source
expression. In this example $D$ contains no free occurrence of its
bound self, so $D[x/s]=D$. Nevertheless the rule is present in the
derivation and has a visible translation.

\textbf{Translate introduction and opening.}
The target literal has precise type $\ty{Lit}$, while the translation
of the source literal's assigned type is $\ty{Exact}$. The object
translation therefore inserts the coercion already constructed:
\[
 \tr{d_{\mathrm{obj}}}
       =\cast{\{A=W\}}{\ev{expose}}.
\]
For recursive opening, write
$P_E=[W\le s.A,\ s.A\le W]$ for the telescope of $\ty{Exact}$ and define
the exact wrapper abbreviation
\[
 R(c)=\ev{foldSelf}_{P_E}(\ev{unfoldSelf}(c)).
\]
The unfold wrapper opens the recursive presentation at the atom's root;
the fold wrapper packages the opened declaration's telescope as an
object interface. These are atom constructors, not evidence composition.
Here opening leaves the telescope's external types unchanged, so
$R(c):\ty{Exact}$ whenever $c:\ty{Exact}$.
Both constructors preserve the root. A view records evaluated interface
facts, as the next calculation will explain; their view rules give
$\view{R(c)}=\view{c}$ when interpreted over a typed store.
They are still present in the syntax and consume normalizer fuel.

Thus the bound derivations for this literal-origin source program
translate to $\ev{lower}_A(R(x))$ and $\ev{upper}_A(R(x))$.
The earlier parameter calculation used $\ev{lower}_A(x)$ and
$\ev{upper}_A(x)$ because it assumed $x:D$ directly.
This is the exact difference between the two derivations.

\textbf{Assemble the whole target.}
The identity lambda translates to itself: its body uses variable typing.
Each of the four lets uses the let clause from
Section~\ref{sec:evidence}. With the chosen lower-bound, upper-bound,
and argument subsumptions, the complete translation is
\[
\tgt{\begin{aligned}
 \ev{let}\ x&=\cast{\{A=W\}}{\ev{expose}}\ \ev{in}\\
 \ev{let}\ f&=\lambda(u:\top).u\ \ev{in}\\
 \ev{let}\ \ty{hidden}&=\cast{f}{\ev{lower}_A(R(x))}\ \ev{in}\\
 \ev{let}\ \ty{recovered}
   &=\cast{\ty{hidden}}{\ev{upper}_A(R(x))}\ \ev{in}\\
 &\ty{recovered}\ \cast{f}{\ev{top}(W)}.
\end{aligned}}
\]
Every constructor introduced by this derivation is now accounted for:
the object rule inserts $\ev{expose}$, \textsc{Rec-E} inserts $R$,
the three subsumptions insert casts, and lambda and let translation
retain the program structure. The notation still suppresses the
implementation's variable indices; $R$ and the member-read abbreviations
have explicit definitions.

\textbf{Close the body by allocation and substitution.}
After the first two bindings have allocated their bare values, the store
contains
\[
 \sigma(\ell)=\{A=W\},\qquad
 \sigma(z)=\lambda(u:\top).u,
 \qquad \Gamma_\sigma(\ell)=\ty{Lit},\quad
        \Gamma_\sigma(z)=W.
\]
The literal's wrapper is retained in the atom substituted for $x$:
\[
 a=\cast{\ell}{\ev{expose}},\qquad a^\circ=R(a),\qquad
 \theta=[x:=a,\ f:=z].
\]
For the remaining body $B$ of the displayed target, substitution gives
\[
\tgt{\begin{aligned}
 B\theta={}&
 \ev{let}\ \ty{hidden}=\cast{z}{\ev{lower}_A(a^\circ)}\ \ev{in}\\
 &\ev{let}\ \ty{recovered}
     =\cast{\ty{hidden}}{\ev{upper}_A(a^\circ)}\ \ev{in}\\
 &\ty{recovered}\ \cast{z}{\ev{top}(W)}.
\end{aligned}}
\]
In evidence receivers, $R(x)$ has become the whole atom $R(a)$.
In type selections, $x.A$ has become $\ell.A$, since
$\ev{root}(a)=\ell$. Substituting the two aliases next gives the
closed application
\[
 \cast{\cast{z}{\ev{lower}_A(a^\circ)}}{\ev{upper}_A(a^\circ)}
 \quad \cast{z}{\ev{top}(W)}.
\]
This is the application evaluated in Section~\ref{sec:application}.
The following view calculation supplies its remaining obligation:
the two member reads really normalize to identity at this store.
The companion's \texttt{FullSource} examples check the full source
derivation and its exact target translation; the machine calculation here
uses the stated allocation and substitution rules.
